% ============================================================
%  第一部分：案例设定、收益最大化问题与逐步数值求解
% ============================================================

\section{案例设定：双产出技术与 Shephard 转换函数}
\label{sec:convention}

\subsection{一个"厨房"里的两产品企业}

设想一家面包房，给定一天内可支配的全部要素（一名师傅、一台烤箱、固定数量的面粉与糖），
它可以生产两种产品：

\begin{itemize}
  \item $y_1$：面包的产量（单位：件）；
  \item $y_2$：蛋糕的产量（单位：件）。
\end{itemize}

两种产品都要占用同一批要素，因此之间存在替代关系：多烤一个蛋糕，
就得少烤几个面包。这种"给定要素禀赋下两种产出的技术替代关系"，
在数学上不必回到生产函数去处理，而可以直接用一个\textbf{转换函数}
（transformation function，Shephard 的标准设定）来描述。
本案例的技术设定为
\begin{equation}
F(y_1,y_2)=y_1+\frac{y_2^{2}}{4}-10,
\label{eq:tech}
\end{equation}
市场价格为 $p_1=2$、$p_2=4$（单位：元/件）。目标是选择产出组合最大化收益：
\begin{equation}
\max_{y_1,y_2}\ \pi=p_1y_1+p_2y_2=2y_1+4y_2
\quad\text{s.t.}\quad F(y_1,y_2)=0.
\label{eq:rm}
\end{equation}

在进一步讨论之前，必须先固定 $F$ 的\textbf{符号约定}。
这是整件事的根基：约定一旦确定，$- \lambda F$ 还是 $+\lambda F$、
$\lambda$ 的正负、可行集在哪一侧，就全部锁定了。

\subsection{转换函数的三项约定}

\begin{definition}[转换函数及其约定]
设 $F:\R^{2}_{+}\to\R$ 连续可微，且对每个 $i$ 有 $\partial F/\partial y_i>0$
（产出越多，$F$ 越大，即技术约束是"紧"的、不存在自由处置意义上的浪费被隐藏）。
称 $F$ 为该技术的\textbf{转换函数}，并约定：
\begin{equation}
F(y_1,y_2)
\begin{cases}
=0, & \text{产出组合恰好落在生产可能性边界（PPF）上，技术有效；}\\[2pt]
<0, & \text{产出组合落在 PPF 内部，资源未用尽，技术无效率；}\\[2pt]
>0, & \text{产出组合超出 PPF，技术上不可行。}
\end{cases}
\label{eq:convention}
\end{equation}
\end{definition}

三条约定的含义要逐条落实。

\begin{itemize}
  \item \textbf{为什么 $F=0$ 是边界。}等值线 $F(y)=0$ 是一条曲线，
        它把平面分成两块区域；由于 $F$ 连续，从 $F>0$ 的点走到 $F<0$ 的点
        必定穿过这条曲线，所以它就是集合 $\{y:F(y)\le 0\}$ 的边界，
        也就是生产可能集的上边界。
  \item \textbf{为什么可行的一侧是 $F\le 0$ 而不是 $F\ge 0$。}
        关键在于单调性假设 $\partial F/\partial y_i>0$：
        若某点满足 $F<0$，则把任一产出减少一点点，$F$ 会变得更小，
        于是该点连同它"左下方的所有点"都属于可行集；
        而若 $F>0$，减少产出只能使 $F$ 沿着连续路径下降，
        无论减多少都不可能把它变进 $F\le 0$ 而不丢掉已实现的一部分产出，
        因此 $F>0$ 的点位于要素禀赋所支持的范围之外。
        一句话：$\partial F/\partial y_i>0$ 保证了"内部点可达、外部点不可达"。
  \item \textbf{为什么必须是严格不等号才叫无效率。}
        $F<0$ 表示存在未用尽的要素：可以在不减少 $y_1$ 的条件下增加一些 $y_2$
        （或反之），直到 $F$ 上升到 $0$。因此利润最大化企业绝不会停在严格内部点上。
\end{itemize}

\begin{remark}[约定是任意的，但不能中途更换]
把可行的一侧写成 $F\le 0$ 只是约定。若有人用\ $\widetilde{F}=-F$，
则同一技术的可行集变为 $\widetilde{F}\ge 0$，拉格朗日要写成
$\mathcal{L}=p\cdot y+\mu\,\widetilde{F}$ 才能得到正的乘子。
\textbf{判断一套写法是否自洽，只有一条标准：在最优点上，价格向量必须与所选转换函数的梯度同向且成严格正比例。}
本案例始终遵守定义 \ref{eq:convention}，故拉格朗日写为减去 $\lambda F$，乘子为正。
\end{remark}

除上述三项约定之外还有一层容易被忽略的问题：$F$ 本身的量纲。\eqref{eq:tech} 中的 $\frac14 y_2^{2}$ 看似与 $y_1$ 量纲不合，
严格写法应引入带单位的曲率参数 $\kappa$：
\begin{equation}
F(y_1,y_2)=y_1+\frac{\kappa}{2}y_2^{2}-c,
\qquad [\kappa]=\text{件}^{-1},\quad [c]=\text{件},
\label{eq:tech-general}
\end{equation}
本案例取 $\kappa=\tfrac12$、$c=10$，于是 $\frac{\kappa}{2}y_2^{2}=\frac14 y_2^{2}$，
恰为 \eqref{eq:tech}。此时 $[F]=\text{件}$，
$F_1=1$ 无量纲，$F_2=\kappa y_2$ 也无量纲。
后面 \S\ref{sec:grad} 会用到这一点：$\lambda$ 的单位是元/件，与价格一致。

\subsection{三项约定的数值检验}

取同一条竖线 $y_2=4$ 上的三个点来检验约定 \eqref{eq:convention}：

\begin{equation}
\begin{array}{lll}
\text{点 }(y_1,y_2)=(6,4): & F=6+\tfrac{16}{4}-10=0, & \text{前沿（技术有效）；}\\[2pt]
\text{点 }(y_1,y_2)=(4,4): & F=4+\tfrac{16}{4}-10=-2<0, & \text{内部（要素闲置）；}\\[2pt]
\text{点 }(y_1,y_2)=(8,4): & F=8+\tfrac{16}{4}-10=+2>0, & \text{不可行。}
\end{array}
\label{eq:check-three}
\end{equation}

\begin{figure}[H]
\centering
\begin{tikzpicture}[x=0.95cm,y=0.50cm,line cap=round,line join=round]
  % ---- 可行集（F<=0）阴影 ----
  \fill[black!6] (0,0) -- (0,10)
    -- plot[domain=0:6.3246,samples=70] (\x,{10-\x*\x/4})
    -- (6.3246,0) -- cycle;
  % ---- 坐标轴 ----
  \draw[axisline] (0,0) -- (7.6,0) node[below,axislabel] {$y_2$（蛋糕，件）};
  \draw[axisline] (0,0) -- (0,13.8) node[left,axislabel] {$y_1$（面包，件）};
  % ---- 生产可能性边界 ----
  \draw[curveA,domain=0:6.3246,samples=90] plot (\x,{10-\x*\x/4});
  \node[figlabel,anchor=east] at (6.0,0.55) {$y_1=10-\dfrac{y_2^{2}}{4}$};
  % ---- 检验点所在的竖线 ----
  \draw[curveC] (4,0) -- (4,8.6);
  % ---- 三个检验点 ----
  \draw[nodewhite] (4,6) circle (2.2pt);
  \draw[nodewhite] (4,4) circle (2.2pt);
  \draw[nodewhite] (4,8) circle (2.2pt);
  \fill (4,6) circle (2.2pt);
  \fill (4,4) circle (2.2pt);
  \fill (4,8) circle (2.2pt);
  % ---- 三个标签 ----
  \node[figlabel,anchor=west] at (4.25,8.0) {$(8,4)$：$F=+2>0$，技术不可行};
  \node[figlabel,anchor=west] at (4.25,6.0) {$(6,4)$：$F=0$，生产前沿};
  \node[figlabel,anchor=west] at (4.25,4.0) {$(4,4)$：$F=-2<0$，内部无效率};
  % ---- 区域说明 ----
  \node[figlabel] at (2.2,2.2) {$F<0$：可行域};
  \node[figlabel] at (1.6,11.6) {$F>0$：不可行域};
  % ---- 端点 ----
  \draw (0,0.12) -- (-0.12,0.12) -- (-0.12,0);
  \node[figlabel,anchor=east] at (0.05,10) {$10$};
  \node[figlabel,anchor=north] at (6.3246,0.05) {$\sqrt{40}$};
\end{tikzpicture}
\caption{转换函数的符号约定：粗实线为生产可能性边界 $F(y_1,y_2)=0$，
  其左下方的浅灰区域为可行集 $\{y:F(y)\le 0\}$，右上区域技术不可行。
  同一竖线 $y_2=4$ 上的三个点分别对应 $F>0$、$F=0$、$F<0$ 三种情形。}
\label{fig:convention}
\end{figure}
\figexplan{横轴为 $y_2$（蛋糕），纵轴为 $y_1$（面包）。
粗实线 $y_1=10-y_2^{2}/4$ 是唯一的技术前沿，它向下弯曲，
反映 $y_2$ 越多、再多烤一单位蛋糕要放弃的面包越多（机会成本递增）。
浅灰阴影为可行集，其与两坐标轴之间的部分就是"要素未用尽"的内部点。
虚线竖线穿过三个检验点：最低点 $(4,4)$ 在阴影内（$F=-2$），
中间点 $(6,4)$ 恰在边界上（$F=0$），最高点 $(8,4)$ 在阴影外（$F=+2$），
它在给定要素下无法生产出来。}

\subsection{本案例技术的三条性质}

在动手求解之前，先把技术的形状说清楚，因为二阶条件与唯一性都取决于它。

\begin{enumerate}
  \item \textbf{单调性}。$F_1=\partial F/\partial y_1=1>0$，
        $F_2=\partial F/\partial y_2=\dfrac{y_2}{2}>0$（对 $y_2>0$），
        满足转换函数的单调性要求，故约定 \eqref{eq:convention} 成立。
  \item \textbf{可行集紧（闭且有界）}。由 $F\le 0$ 且 $y\ge 0$ 得
        $0\le y_2\le \sqrt{40}$、$0\le y_1\le 10-\dfrac{y_2^{2}}{4}$，
        可行集是有界闭集；连续函数 $p\cdot y$ 在其上必取到最大值（Weierstrass 定理），
        \textbf{最大值一定存在}，这是后面一切讨论的前提。
  \item \textbf{前沿严格凹、机会成本递增}。把前沿写成显式函数
        \begin{equation}
        y_1=T(y_2)=10-\frac{y_2^{2}}{4},\qquad
        T'(y_2)=-\frac{y_2}{2}<0,\quad T''(y_2)=-\frac12<0,
        \label{eq:ppf-explicit}
        \end{equation}
        $T$ 严格凹，即生产可能集严格凸之外的"边界向外凸"形状：
        随着 $y_2$ 增加，$-(dy_1/dy_2)=y_2/2$ 递增，
        每多烤一单位蛋糕要放弃的面包越来越多。
\end{enumerate}

第三条性质直接给出\textbf{二阶充分条件}：把约束代入目标，
得到只含一个变量的收益函数
\begin{equation}
\pi(y_2)=p_1T(y_2)+p_2y_2
=2\Bigl(10-\frac{y_2^{2}}{4}\Bigr)+4y_2
=20-\frac{y_2^{2}}{2}+4y_2,
\qquad \pi''(y_2)=-1<0,
\label{eq:one-dim}
\end{equation}
它是 $y_2$ 的严格凹函数。因此任何满足一阶条件的内点都是\textbf{唯一的全局最大点}，
不需要再讨论多峰或局部极值问题。这一点在课堂上值得先说，
因为它把"检查二阶条件"变成了看一眼二阶导数符号。

完成平方还能把 \eqref{eq:one-dim} 写成一个更便于口算的形式：
\begin{equation}
\pi(y_2)=28-\frac{(y_2-4)^{2}}{2}.
\label{eq:pi-complete-square}
\end{equation}
沿前沿的收益是一条以 $y_2=4$ 为顶点、最高值为 $28$ 的抛物线。
后面 \S\ref{sec:solve} 用拉格朗日法得到的答案，必须与这条抛物线的顶点完全一致。

\section{收益最大化问题与拉格朗日一阶条件}
\label{sec:foc}

\subsection{问题与几何图像}

把 \eqref{eq:rm} 写成一个正式的最优化问题：
\begin{equation}
\tag{RM}
\max_{y_1,y_2}\ p_1y_1+p_2y_2
\quad\text{s.t.}\quad F(y_1,y_2)=0,\quad y_1\ge 0,\ y_2\ge 0.
\label{eq:rm-formal}
\end{equation}
目标函数 $p\cdot y$ 是线性的，其等值线$p_1y_1+p_2y_2=\bar{\pi}$ 是直线，
斜率（以横轴 $y_2$、纵轴 $y_1$ 度量）为
\begin{equation}
\frac{dy_1}{dy_2}=-\frac{p_2}{p_1}=-\frac{4}{2}=-2.
\label{eq:isor-slope}
\end{equation}
几何上，\eqref{eq:rm-formal} 就是问：在给定斜率的直线族中，
哪一条与前沿 $F=0$ 相交且截距（即收益）最大。
答案是与前沿\textbf{相切}的那一条——这就是后面所有代数要复述的内容。

\subsection{拉格朗日函数与一阶条件}

按定义 \ref{eq:convention} 的约定构造拉格朗日函数：
\begin{equation}
\mathcal{L}(y_1,y_2,\lambda)=p_1y_1+p_2y_2-\lambda F(y_1,y_2).
\label{eq:lagrange}
\end{equation}
对两个产出求偏导并令其为零：
\begin{equation}
\begin{cases}
\dfrac{\partial \mathcal{L}}{\partial y_1}=p_1-\lambda\dfrac{\partial F}{\partial y_1}=0,\\[8pt]
\dfrac{\partial \mathcal{L}}{\partial y_2}=p_2-\lambda\dfrac{\partial F}{\partial y_2}=0,
\end{cases}
\quad\Longleftrightarrow\quad
\begin{cases}
p_1=\lambda F_1,\\
p_2=\lambda F_2.
\end{cases}
\label{eq:foc}
\end{equation}
两式相除消去 $\lambda$（要求 $F_2\neq 0$），得到本案例的核心等式：
\begin{equation}
\boxed{\ \frac{p_1}{p_2}=\frac{F_1}{F_2}\ }
\label{eq:core}
\end{equation}
第三个一阶条件是把 $\mathcal{L}$ 对 $\lambda$ 求偏导，它把约束还回来：
$\partial \mathcal{L}/\partial \lambda=-F(y_1,y_2)=0$，即 $F=0$。

三个未知数 $y_1,y_2,\lambda$，三个方程，可解。
这里只需要一条约束规格：$\nabla F=(F_1,F_2)\neq (0,0)$，
在最优点附近由 $F_1=1$ 自动满足（梯度秩为 $1$）。

\subsection{为什么乘子为正：符号、量纲与经济含义}

由 \eqref{eq:foc} 直接读出：
\begin{equation}
\lambda=\frac{p_1}{F_1}=\frac{p_2}{F_2}.
\label{eq:lambda-ratio}
\end{equation}
本例中 $p_1,p_2>0$ 且 $F_1,F_2>0$，故 $\lambda>0$。这不是巧合：

\begin{itemize}
  \item \textbf{符号方面}。$F$ 沿任一产出增加的方向上升，而价格非负，
        两者之比自然为正。拉格朗日里写 $-\lambda F$ 而不是 $+\lambda F$，
        正是为了保证这个正数出现；若写成 $+\lambda F$，
        解出的 $\lambda$ 会是 $-2$，解释为"价格与技术梯度反向"，徒增混乱。
  \item \textbf{量纲方面}。由 \eqref{eq:tech-general}，
        $[F_1]=1$（无量纲），$[p_1]=$ 元/件，故 $[\lambda]=$ 元/件，
        与价格同量纲。于是 $p=\lambda\nabla F$ 两边量纲一致：
        无单位的梯度方向被乘以一个"元/件"的比例因子，得到价格向量。
  \item \textbf{经济方面}。$\lambda$ 把\textbf{纯技术的方向}
        （$\nabla F$，不含任何价格信息）缩放成\textbf{市场价格的方向}。
        \S\ref{sec:shadow} 将证明它等于收益对技术水平的导数，
        是技术约束的影子值。
\end{itemize}

\section{逐步求解：从梯度到最优产出组合}
\label{sec:solve}

下面按四个步骤把 \eqref{eq:rm-formal} 解到底，每一步的中间结果都单独列出，
便于课堂上逐步板书。

\subsection{第一步：求偏导（梯度分量）}

由 $F=y_1+\dfrac{y_2^{2}}{4}-10$：
\begin{equation}
F_1=\frac{\partial F}{\partial y_1}=1,
\qquad
F_2=\frac{\partial F}{\partial y_2}=\frac{y_2}{2}.
\label{eq:grad-components}
\end{equation}
注意 $F_1$ 是常数，$F_2$ 随 $y_2$ 递增：这条技术对 $y_1$ 而言是"线性的"，
对 $y_2$ 而言呈凸的技术代价，这正是前沿弯曲的来源。

\subsection{第二步：代入一阶条件}

把 \eqref{eq:grad-components} 与价格 $p_1=2,p_2=4$ 代入 \eqref{eq:foc}：
\begin{equation}
2=\lambda\cdot 1,
\qquad
4=\lambda\cdot \frac{y_2}{2}.
\label{eq:foc-numbers}
\end{equation}
第一式立刻给出
\begin{equation}
\lambda=2.
\label{eq:lambda-value}
\end{equation}
代入第二式：$4=2\cdot \dfrac{y_2}{2}=y_2$，故
\begin{equation}
y_2^{*}=4.
\label{eq:y2-star}
\end{equation}

\subsection{第三步：回代前沿求出产出组合与收益}

把 $y_2^{*}=4$ 代入前沿 $F=0$，即显式形式 \eqref{eq:ppf-explicit}：
\begin{equation}
y_1^{*}=10-\frac{(4)^{2}}{4}=10-4=6.
\label{eq:y1-star}
\end{equation}
对应的最优收益（最大可能的总收入）为
\begin{equation}
\pi^{*}=p_1y_1^{*}+p_2y_2^{*}=2\times 6+4\times 4=12+16=28.
\label{eq:pi-star}
\end{equation}
最优解汇总：
\begin{equation}
\boxed{\ (y_1^{*},y_2^{*})=(6,4),\qquad \lambda^{*}=2,\qquad \pi^{*}=28\ }
\label{eq:solution}
\end{equation}

\subsection{第四步：验证核心等式}

核心等式 \eqref{eq:core} 的左边是市场价格比，右边取在最优点 $(6,4)$ 处的技术梯度比：
\begin{equation}
\frac{p_1}{p_2}=\frac{2}{4}=\frac12,
\qquad
\frac{F_1}{F_2}\bigg|_{(6,4)}=\frac{1}{4/2}=\frac{1}{2}.
\label{eq:core-check}
\end{equation}
两边相等，验证通过。同时检查 $p=\lambda\nabla F$：
\begin{equation}
\lambda\nabla F=2\cdot\Bigl(1,\frac{y_2}{2}\Bigr)\bigg|_{y_2=4}
=2\cdot(1,2)=(2,4)=(p_1,p_2)=p.
\label{eq:p-lambda-grad}
\end{equation}
价格向量与技术梯度在最优点上共线，比例因子为 $\lambda=2$。

\subsection{交叉验证一：消元法}

不做拉格朗日，改用"把约束消掉"的老办法。由前沿 $y_1=10-\dfrac{y_2^{2}}{4}$
代入目标，得到单变量问题
\begin{equation}
\max_{0\le y_2\le \sqrt{40}}\ \pi(y_2)=20-\frac{y_2^{2}}{2}+4y_2
=28-\frac{(y_2-4)^{2}}{2}.
\label{eq:onedim}
\end{equation}
一阶条件：$\pi'(y_2)=4-y_2=0\Rightarrow y_2^{*}=4$；
二阶条件：$\pi''=-1<0$，严格凹，故为唯一全局最大点。回代得 $y_1^{*}=6$、$\pi^{*}=28$。
与 \eqref{eq:solution} 完全一致。

两种方法的对应关系值得点明：消元法里的"$4-y_2$"就是"$p_2-\lambda F_2$"
除以某个正因子的结果，二者是同一条边际条件的两种写法；
拉格朗日法的优势在于不必显式解出前沿函数 $T(\cdot)$，
在多产出、多约束的问题里尤其明显。

\subsection{交叉验证二：前沿上的收益对照表}

沿前沿逐点计算收益，可以直接看到最大值出现在何处。
由 \eqref{eq:pi-complete-square}，收益只取决于 $(y_2-4)^{2}$，
表 \ref{tab:frontier-pi} 给出七个代表性组合。

\begin{table}[H]
\centering
\caption{沿生产可能性边界的收益对照（$p_1=2,\ p_2=4$）}
\label{tab:frontier-pi}
\begin{tabular}{ccccc}
\toprule
$y_2$  & $y_1=10-\dfrac{y_2^{2}}{4}$ & $\pi=2y_1+4y_2$ & $28-\dfrac{(y_2-4)^{2}}{2}$ & 位置 \\
\midrule
$0$      & $10.00$  & $20.00$  & $20.00$  & 端点（只产面包）\\
$2$      & $9.00$   & $26.00$  & $26.00$  & 偏离 \\
$3$      & $7.75$   & $27.50$  & $27.50$  & 偏离 \\
$\mathbf{4}$ & $\mathbf{6.00}$ & $\mathbf{28.00}$ & $\mathbf{28.00}$ & \textbf{最优（相切）}\\
$5$      & $3.75$   & $27.50$  & $27.50$  & 偏离 \\
$6$      & $1.00$   & $26.00$  & $26.00$  & 偏离 \\
$\sqrt{40}\approx 6.325$ & $0$ & $25.30$ & $25.30$ & 端点（只产蛋糕）\\
\bottomrule
\end{tabular}
\end{table}

表 \ref{tab:frontier-pi} 说明两件事。第一，最大值确实出现在 $y_2=4$ 而不是别处；
第二，两个端点都明显吃亏：只产面包得到 $20$，只产蛋糕得到 $25.30$，
都低于混合生产的 $28$。也就是说，\textbf{只要价格比不极端，
专业化（走到端点）就不是最优}，这一点在 \S\ref{sec:corner} 会被严格化。

\begin{figure}[H]
\centering
\begin{tikzpicture}[x=0.95cm,y=0.50cm,line cap=round,line join=round]
  % ---- 可行集阴影 ----
  \fill[black!6] (0,0) -- (0,10)
    -- plot[domain=0:6.3246,samples=70] (\x,{10-\x*\x/4})
    -- (6.3246,0) -- cycle;
  % ---- 坐标轴 ----
  \draw[axisline] (0,0) -- (7.6,0) node[below,axislabel] {$y_2$（蛋糕，件）};
  \draw[axisline] (0,0) -- (0,13.8) node[left,axislabel] {$y_1$（面包，件）};
  % ---- 生产前沿 ----
  \draw[curveA,domain=0:6.3246,samples=90] plot (\x,{10-\x*\x/4});
  \node[figlabel,anchor=north east] at (1.15,9.68) {PPF：$F=0$};
  % ---- 等收益线 pi = 20（过端点 (10,0)）----
  \draw[curveB] (0,10) -- (5,0);
  \node[figlabel,anchor=north west] at (0.85,8.3) {$\bar\pi=20$};
  % ---- 等收益线 pi = 25.30（过端点 (0,sqrt40)）----
  \draw[curveC] (0,12.649) -- (6.3246,0);
  \node[figlabel,anchor=south west] at (0.15,12.3) {$\bar\pi\approx 25.30$};
  % ---- 等收益线 pi = 28（与 PPF 相切）----
  \draw[thick,black] (2,10) -- (7,0);
  \node[figlabel,anchor=south west] at (6.0,1.6) {$\bar\pi=28$};
  % ---- 最优 -----端点标记
  \fill (4,6) circle (2.6pt);
  \node[figlabel,anchor=south east] at (3.9,6.15) {$(y_2^{*},y_1^{*})=(4,6)$};
  \draw (0,0.12) -- (-0.12,0.12) -- (-0.12,0);
  \node[figlabel,anchor=east] at (0.05,10) {$10$};
  \node[figlabel,anchor=north] at (6.3246,0.05) {$\sqrt{40}$};
\end{tikzpicture}
\caption{收益最大化的几何：粗实线为生产前沿，
  三条直线为不同收益水平的等收益线。$2y_1+4y_2=28$ 与前沿相切于 $(y_2,y_1)=(4,6)$，
  是所有与前沿有交点的等收益线中位置最高的一条。}
\label{fig:optimal-tangency}
\end{figure}
\figexplan{横轴 $y_2$、纵轴 $y_1$。
粗实线（PPF）向下弯曲；三条等收益线斜率同为 $-p_2/p_1=-2$：
虚线 $\bar\pi=20$ 经过只产面包的端点，
点线 $\bar\pi\approx 25.30$ 经过只产蛋糕的端点，
实线 $\bar\pi=28$ 与前沿恰好相切，黑点即最优点 $(y_2^{*},y_1^{*})=(4,6)$。
在穿过相同两条直线的意义上，$\bar\pi=28$ 的位置"最高"：
任何更高的平行线都会与前沿脱离接触，对应技术做不到的产出组合。}
